\chapter{The Exam, the Marks and the Certificate}\label{ch:exam}

\emph{How is the Level 1 exam run, and what does it take to pass?}
The candidate must know the two stages of the exam, the marks, the pass rule, what happens after a failed stage, and the certificate rules. These points follow the Board's \emph{Guideline for Yoga Professionals}, version 0.3, and the Board syllabus for Yoga Protocol Instructor, version 0.2.

\begin{caution}
Rules, fees and dates must be re-read from the Board before use. The Board issues site notices that change some validity dates, and amounts change. Do not rely on any fee or date in this book without checking the current guideline and notices, or asking the institution. The official guidebook was written for an older scheme and must not be used for the exam rules.
\end{caution}

\section*{The level}

This is the first level of the Board's scheme, the Yoga Protocol Instructor. The holder teaches the basics of the common yoga protocol for prevention of disease and promotion of health, in settings such as parks and housing societies. The higher levels (Yoga Wellness Instructor, Yoga Teacher and Evaluator) are out of scope here.

\section*{Who can sit, and training time}

\begin{itemize}
\item Open candidates need no qualification.
\item There is no age limit.
\item Completion of secondary school is suggested for admission to a course, but a training institution may set its own rule.
\item The course is not less than 200 hours, or not less than 3 months part time, or not less than 1 month full time.
\item The certificate carries 12 credit points.
\end{itemize}

\section*{Two stages}

Both stages are compulsory: a theory exam and a practical assessment. Absence in either stage counts as a fail. The assessing institution may, in an exceptional case and on merit, allow a candidate who sat the theory to take the practical on another date.

The theory stage tests what the candidate knows. The practical stage tests whether the candidate can do the practices and teach them.

\section*{Marks}

\begin{center}
\begin{tabular}{@{}lr@{}}
\toprule
Part & Marks \\
\midrule
Theory, split equally over three units & 60 \\
Demonstration skills & 80 \\
Teaching skills & 40 \\
Application of knowledge & 10 \\
Field experience & 10 \\
Practical, total & 140 \\
\midrule
All together & 200 \\
\bottomrule
\end{tabular}
\end{center}

\section*{The theory paper}

\begin{itemize}
\item 60 multiple-choice questions of 1 mark each, in two hours.
\item Answered on an optical mark sheet.
\item No negative marking.
\item Questions are spread equally over the units.
\item The paper is bilingual, in English and Hindi. If the two versions differ, the English one is correct.
\item The candidate returns the question paper and the answer sheet.
\end{itemize}

\section*{The practical assessment}

\begin{itemize}
\item Each candidate is assessed for not more than one hour. The assessor decides the actual length within that limit.
\item The candidate shows a report with photographs of a training programme conducted, and a feedback form. See the Field Experience chapter (\ref{ch:field-experience}).
\end{itemize}

\section*{Pass rule}

\begin{itemize}
\item At least 70 percent in theory and at least 70 percent in practical, each on its own.
\item Up to 5 percent grace marks may be given in one of the two stages.
\item A decimal result is rounded up.
\end{itemize}

Treat the pass mark per stage as the real target. A strong theory score does not rescue a weak practical.

\section*{If a stage is failed}

\begin{itemize}
\item Practical passed, theory failed: the candidate may sit the theory once more within 3 months of the result, paying the prescribed fee.
\item Theory passed, practical failed: the candidate counts as failed and applies again as a fresh candidate, paying the full assessment fee but not a new enrolment fee. The enrolment number stays valid.
\end{itemize}

\section*{Re-checking}

\begin{itemize}
\item A candidate who disagrees with a theory result may ask the institution to re-check the answer sheet, in the candidate's presence, within 15 days of the result, for a fee.
\item The agency replies within 30 days.
\item Only the arithmetic is re-checked.
\item The candidate may not keep a copy of the paper or the answer sheet.
\end{itemize}

\section*{The certificate}

\begin{itemize}
\item Issued by the Board with a unique number and a security code. A digital copy can be downloaded after login.
\item For this level the validity is 5 years, counted from the day the result is declared.
\item The holder renews by attending the Continuing Yoga Education Programme before the certificate expires.
\end{itemize}

\section*{Money}

\begin{itemize}
\item The enrolment fee is paid once for a type of certification and is valid for life.
\item The examination fee is fixed by the institution that holds the exam.
\item There are separate fees for re-checking and for a duplicate certificate.
\item Amounts change. Read them from the Board or the institution before use.
\end{itemize}

\section*{On the day}

\begin{itemize}
\item The admit card is issued at least a week before and can be downloaded from the Board's site.
\item Carry it with a valid identity card and a pen.
\item Reach the centre preferably an hour early.
\item A candidate more than 15 minutes late is not allowed in.
\item Results are declared within 30 days.
\end{itemize}

\section*{What past candidates report}

This is not from the Board and may vary by agency and year. Use it as practice only, and follow the syllabus.
\begin{itemize}
\item The practical takes about an hour, online, in a small batch with a few examiners, with one-to-one questions.
\item Keep the enrolment number and an identity card ready.
\item Place the mat sideways so the whole body is visible on camera.
\item Have a good connection and a clear voice.
\item Answer only what is asked. Say plainly that you do not know rather than guess.
\item Stay calm.
\end{itemize}

\begin{recall}
\begin{itemize}
\item Two stages, both compulsory. Absence counts as a fail.
\item Marks: 200 in all. Theory 60, practical 140.
\item Practical: demonstration 80, teaching 40, application of knowledge 10, field experience 10.
\item Theory: 1 mark each, optical mark sheet, no negative marking. English is correct if versions differ.
\item Pass: 70 percent in each stage on its own. Up to 5 percent grace in one stage.
\item Theory failed: sit again within 3 months. Practical failed: apply again as a fresh candidate.
\item Certificate valid 5 years from the day the result is declared. Renew with the Continuing Yoga Education Programme.
\item Re-read rules, fees and dates from the Board before use.
\end{itemize}
\end{recall}

\subsection*{Worked examples}

\textbf{Describe the two stages of the Level 1 exam and the pass rule.} Two stages, both compulsory: a theory exam and a practical assessment. Absence in either counts as a fail. Marks are 200 in all: theory 60, practical 140. The practical is split as demonstration skills 80, teaching skills 40, application of knowledge 10 and field experience 10. Pass rule: at least 70 percent in theory and at least 70 percent in practical, each on its own. Up to 5 percent grace marks may be given in one of the two stages, and a decimal result is rounded up.

\textbf{Working out the pass marks.} Theory is out of 60, so 70 percent is 42 marks. Practical is out of 140, so 70 percent is 98 marks.
\begin{itemize}
\item Find 70 percent of each stage separately.
\item Compare each score with its own pass mark. A high score in one stage does not make up for a low score in the other.
\item Read the next step from the stage that was failed.
\end{itemize}
A candidate with 45 in theory and 100 in practical has passed both stages. A candidate with 45 in theory and 90 in practical has passed theory but not practical, so the candidate counts as failed and applies again as a fresh candidate, paying the full assessment fee but not a new enrolment fee. A candidate with 38 in theory and 100 in practical has passed the practical and not the theory, and may sit the theory once more within 3 months of the result. How the grace marks are counted is not spelled out in the guideline, so ask the institution.

\textbf{After a failed stage.} A candidate passes the practical but scores below the theory pass mark. The candidate may sit the theory once more within 3 months of the result, paying the prescribed fee. The enrolment fee is not paid again, because it is paid once for a type of certification and is valid for life. In the opposite case the candidate applies again as a fresh candidate, and the enrolment number stays valid.

\subsection*{Common questions}

\begin{itemize}
\item \emph{Do I need to finish school to sit the exam?} Open candidates need no qualification and there is no age limit. Completion of secondary school is only suggested for admission to a course, and a training institution may set its own rule.
\item \emph{How long must the training be?} Not less than 200 hours, or not less than 3 months part time, or not less than 1 month full time.
\item \emph{Does the enrolment fee have to be paid again if I fail?} No. A candidate who fails the practical pays the full assessment fee again but not a new enrolment fee.
\item \emph{Can I use the official guidebook for the exam rules?} No. It was written for an older scheme. Use the Board's current guideline and notices.
\item \emph{Where do the application-of-knowledge and field-experience marks come from?} They are part of the practical assessment. The guideline gives the marks but not the detailed content. Ask the institution what it assesses.
\item \emph{Can I trust the fees and dates in this book?} Not fully. Read the current guideline and notices, or ask the institution, before relying on any fee or date.
\end{itemize}

\subsection*{Check yourself}

\begin{enumerate}
\item What is the pass mark in the theory stage?
\item A candidate passes the practical and fails the theory. What can the candidate do?
\item Is there negative marking in the theory paper?
\item How long is a Level 1 certificate valid, and how is it renewed?
\item If the English and Hindi versions of a question differ, which one is correct?
\item What is the total number of marks, and how are they split between theory and practical?
\item Fill in the blank: the theory paper has multiple-choice questions of \underline{\hspace{1cm}} mark each, answered on an optical mark sheet, with \underline{\hspace{1cm}} negative marking.
\item True or false: a Level 1 certificate is valid for 3 years.
\item Match the practical component with its marks: demonstration skills, teaching skills, application of knowledge, field experience (marks: 10, 10, 40, 80).
\item A candidate scores 40 in theory and 110 in practical. Has the candidate passed? What is the next step?
\item A candidate who sat the theory cannot attend the practical on the set date. What does the guideline say?
\item A candidate disagrees with the theory result. What can the candidate ask for, and what are the limits?
\item Without looking, write the minimum marks needed in theory and in practical. Can a high score in one make up for a low score in the other?
\item Draw a small decision tree: the result in each stage (pass or fail) and the next step for each of the four combinations.
\end{enumerate}

\subsection*{Answers}

\begin{enumerate}
\item 70 percent of the theory marks, separate from the practical.
\item Sit the theory once more within 3 months of the result, after paying the prescribed fee.
\item No.
\item Five years, counted from the day the result is declared. Renew by attending the Continuing Yoga Education Programme before it expires.
\item The English one.
\item 200 in all: theory 60 and practical 140.
\item 1 mark each; no negative marking.
\item False. It is valid for 5 years.
\item Demonstration skills 80, teaching skills 40, application of knowledge 10, field experience 10.
\item The theory pass mark is 42, so the theory is not passed. The practical pass mark is 98, so the practical is passed. The candidate may sit the theory once more within 3 months of the result, paying the prescribed fee.
\item Absence counts as a fail, except that the assessing institution may, in an exceptional case and on merit, allow the candidate to take the practical on another date.
\item The candidate may ask the institution to re-check the answer sheet, in the candidate's presence, within 15 days of the result, for a fee. The agency replies within 30 days. Only the arithmetic is re-checked, and the candidate may not keep a copy of the paper or answer sheet.
\item Theory 42 out of 60 and practical 98 out of 140. No, each stage is passed on its own.
\item Both passed: pass. Practical passed and theory failed: sit the theory once more within 3 months. Theory passed and practical failed: apply again as a fresh candidate. Both failed: the notes do not say, so ask the institution.
\end{enumerate}

\paragraph{Sources.} Yoga Certification Board, \emph{Guideline for Yoga Professionals}, version 0.3 (effective 27 February 2020), and the Board syllabus for Yoga Protocol Instructor, version 0.2; Guideline for the Yoga Protocol examination, version 0.2. The candidate study notes on the practical are practice only, not an authority. All in own words.
